
\section{Related Work}
\label{sec:related}

\noindent \textbf{Vision-Language-Action Models}.
Vision-Language-Action (VLA) models~\citep{kim24openvla,li2023vision,DeeR-VLA,zitkovich2023rt} extend Vision-Language Models (VLMs)~\citep{alayrac2022flamingovisuallanguagemodel,karamcheti2024prismaticvlmsinvestigatingdesign,li2022blipbootstrappinglanguageimagepretraining, liu2023visualinstructiontuning,radford2021learning} by incorporating action generation, bridging the gap between perception and action. These models enable machines to understand visual and textual inputs and generate corresponding actions for tasks~\citep{awadalla2023openflamingoopensourceframeworktraining,li2023vision} such as robotic manipulation and object retrieval. VLA models typically use pretrained VLMs~\citep{karamcheti2024prismaticvlmsinvestigatingdesign} to encode visual and linguistic data into a shared representation, from which actions are generated either as discrete tokens or continuous values. A prominent recent trend within VLA is the adoption of diffusion models for generating coherent continuous action sequences. This paradigm is exemplified by models such as CogACT~\citep{li2024CogACT}, DexVLA~\citep{wen2025dexvla}, DiVLA~\citep{wen2024diffusion}, $\pi_0$~\citep{black2024pi_0}, and TinyVLA~\citep{wen2024tinyvla}. Many of these diffusion-based VLAs employ a componentized design: the foundational VLM processes visual and linguistic inputs to produce a condensed feature representation, which then conditions a distinct diffusion-based action module responsible for the iterative generation of precise action trajectories. This often involves the VLM output steering the denoising process within the specialized action decoder. 


\noindent \textbf{Efficient Vision-Language-Action Models}.
The computational complexity of Vision-Language Models~\citep{liu2025shifting,wen2025stop} poses significant challenges for their real-time deployment, particularly in applications such as robotic control that require rapid decision-making. To address this issue, recent efforts to accelerate VLA models have been primarily categorized into training-aware and training-free methods. Training-aware approaches, such as RoboMamba~\citep{liu2024robomambaefficientvisionlanguageactionmodel}, EfficientVLM~\citep{wang2022efficientvlmfastaccuratevisionlanguage}, and DeeR-VLA~\citep{DeeR-VLA}, focus on optimizing model architectures or applying compression techniques followed by retraining, achieving significant speedups while maintaining performance. For instance, DeeR-VLA reduces computational costs by leveraging dynamic reparameterization and efficient pruning strategies, which enable more flexible and scalable model deployment. Similarly, For example, Mole-VLA~\citep{zhang2025mole} reduces computational costs by dynamically activating only a subset of model layers based on task-specific needs. In contrast, training-free methods, such as VLA-Cache~\citep{xu2025vla}, enhance efficiency by reusing previously computed results for unchanged tokens between consecutive frames, which is particularly beneficial in scenarios with minimal variation in visual input.
